\documentclass[instructions.tex]{subfiles}
\begin{document}
 
\chapter{Des effets de la patience.}
 
Elle adoucit nos peines, elle expie le péché, elle éprouve notre vertu.
 
\primo Les moindres peines, sans la patience, deviennent insupportables. On ne se soulage pas par l'inquiétude et le chagrin, on ne fait qu'aigrir son cœur et irriter le mal. La patience, au contraire, calme notre esprit et adoucit nos inquiétudes.
 
La croix ne nous accable que parce que nous voulons y succomber ; elle ne nous afflige qu'autant que nous voulons nous affliger ; ce qui a fait dire à un ancien que nous n'avons de croix que celles que nous voulons avoir, parce qu'elles ne sont sensibles qu'autant que nous les prenons à cœur. Elles cessent de nous crucifier et de nous affliger, aussitôt que nous les portons avec résignation. Il est vrai que la patience n'ôte pas la peine et la douleur, mais elle en tempère l'amertume ; elle fait, selon la parole du Sauveur, qu'on possède son âme en paix\footnote{In patientiâ vestrâ possidebitis animas vestras. Luc. 21.}.
 
\secundo La patience expie le péché ; en peu de temps elle peut en expier beaucoup. Si nous comprenions ce que nos péchés méritent, loin de nous plaindre des souffrances, nous dirions que Dieu se contente de trop peu de chose ; que, loin de mériter des consolations, nous ne méritons que des châtimens, trop heureux d'échapper à l'enfer à ce prix. Je l'ai mérité cent fois, et je n'y suis pas encore ; je dois donc adorer la miséricorde de mon Dieu, qui, en échange de ces tourmens horribles, se contente de quelques peines en cette vie.
 
N'oublions jamais que, quand nous n'aurions commis qu'un péché mortel, Dieu est d'une si haute majesté, que toutes les peines de cette vie ne sont pas capables de réparer en rigueur l'injure qui lui est faite par un seul crime. Ne nous plaignons donc pas des afflictions dont sa miséricorde nous fait part pour l'expiation de nos péchés, puisqu'elles sont toujours infiniment au-dessous de ce que nous méritons.
 
\tertio Enfin, la patience dans l'adversité éprouve la vertu. Oh ! qu'il est difficile d'allier la prospérité avec la sainteté ! Jouir d'une grande fortune, et servir Dieu dans le détachement ; posséder de grands talens, des emplois éclatans, et aimer l'obscurité ; être chaste, et vivre au milieu des délices ; en un mot, pouvoir tout ce qu'on veut, et ne vouloir que ce qui plaît à Dieu, voilà un rare prodige. Il faut de la grandeur d'âme, dit saint Bernard, pour n'être pas séduit, quand tout nous réussit et que tout nous flatte.
 
Mais c'est une grandeur d'âme encore plus héroïque de conserver la vertu dans les tribulations. Les honneurs, dit-on, changent les mœurs. On parlerait avec plus de vérité, si l'on disait qu'ils les font connaître. Disons la même chose de l'adversité. Celui qui n'a pas été éprouvé, dit le Sage, que sait-il ? peu de choses. Il ne sait servir Dieu qu'imparfaitement, n'ayant point encore donné de grandes marques de son courage. C'est dans l'occasion qu'un soldat qui paie de sa personne et de son sang, fait connaître sa bravoure et sa fidélité. C'est aussi dans les occasions qu'un chrétien fait connaître ce qu'il est.
 
Ce n'est pas une grande merveille d'être patient quand on n'a rien à souffrir, de pratiquer la douceur quand rien ne nous contredit, d'être vertueux quand on n'a point d'occasion d'être méchant ; mais passer sa vie dans la disgrâce, dans l'adversité, dans l'abandon, être accusé, opprimé, exercé par la contradiction et l'imposture, sans être ébranlé, c'est l'héroïsme de la vertu.
 
Où trouver de ces âmes fortes et généreuses qui puissent dire comme saint Paul : Je me plais dans les infirmités, dans la nécessité, dans les persécutions, dans les afflictions pour Jésus-Christ, car c'est alors que je suis plein de force\footnote{2,Cor.13.}. Voilà le partage des grandes âmes ; c'est dans la tribulation que Dieu éprouve et purifie ceux qu'il chérit. C'est parce que vous étiez agréable à Dieu, disait l'ange à Tobie, qu'il était nécessaire que vous fussiez éprouvé. Oh ! que la vertu est faible quand elle n'a pas été éprouvée ! Quelle est rare dans la prospérité et dans le succès !
 
\section{Résolutions}
 
\primo Puisque j'ai besoin de faire pénitence, et que les afflictions qui m'arrivent peuvent me servir à expier mes péchés, je regarderai ces afflictions comme des pénitences légères que Dieu m'envoie pour m'aider à satisfaire à sa justice. --

\secundo Et je ferai tous mes efforts pour ne pas en perdre le fruit. 


\prayer{Mon Dieu, aidez-moi à profiter des épreuves et des peines que j'aurai désormais, pour me purifier devant vous, et payer les dettes immenses que j'ai contractées envers votre justice.}
 
\end{document}
